% GENERATED FILE. Edit canonical lessons, then run npm run generate:books.
% canonical-source-hash: fnv1a64:d45bcfda74e094fb
% canonical-lessons: HI-C403-staind, HI-C403-rahgir, HI-C403-futpath, HI-C403-tahalna, HI-C403-haphna

\chapter{Stands and Footpaths}
\label{ch:hi-stands-and-footpaths}
\hlchaptermodality{403}

\begin{chapteropening}
\textbf{By the end of this chapter:} \emph{I can find the bus stand, walk along the pavement and say I am out of breath.}

Complete the last lesson of chapter 403: I can find the bus stand, walk along the pavement and say I am out of breath.
\end{chapteropening}

\section[\texorpdfstring{\hi{स्टैंड} (sṭaiṇḍ)}{स्टैंड (sṭaiṇḍ)}]{\hi{स्टैंड} (sṭaiṇḍ) — a stand (for buses, taxis or rickshaws)}

\label{lesson:HI-C403-staind}



\begin{quote}
Before the new one: say the Hindi for to doze, then the Hindi for to rest a while.
\end{quote}

\subsection*{You'll want to know: \hi{स्टैंड}}
\textbf{\hi{स्टैंड}} — \emph{sṭaiṇḍ} — \textquotedblleft{}a stand (for buses, taxis or rickshaws)\textquotedblright{}.

\begin{quote}
\hi{बस} \hi{स्टैंड} \hi{यहाँ} \hi{से} \hi{पास} \hi{है।}
\end{quote}

\begin{quote}
\emph{bas sṭaiṇḍ yahā̃ se pās hai.}
\end{quote}

\begin{quote}
The bus stand is close to here.
\end{quote}

Masculine.

\begin{cousinweb}
\textbf{\hi{स्टैंड}} (\emph{sṭaiṇḍ}) is borrowed from English \emph{stand}, the place where buses, taxis or rickshaws wait for passengers: \hi{बस} \hi{स्टैंड}, \hi{रिक्शा} \hi{स्टैंड}.
\end{cousinweb}

\begin{grammarlens}[title={What you've built}]
One more word for getting around.
\end{grammarlens}

\subsection*{Your turn}
\begin{itemize}
\raggedright
  \item \emph{Say it:} \emph{sṭaiṇḍ}
  \item \emph{Say it:} \emph{sṭaiṇḍ}, once more
  \item \emph{Say it:} \emph{bas sṭaiṇḍ yahā̃ se pās hai.} — The bus stand is close to here.
  \item \emph{Recall:} say the Hindi for to doze, then the Hindi for to rest a while, then say \emph{sṭaiṇḍ} again
\end{itemize}

\begin{tcolorbox}[breakable,colback=teal!4,colframe=teal!35!black,title={Before you move on}]
What does \hi{स्टैंड} mean? (\textquotedblleft{}A stand (for buses, taxis or rickshaws)\textquotedblright{}.) Say it once more.
\end{tcolorbox}
\section[\texorpdfstring{\hi{राहगीर} (rāhgīr)}{राहगीर (rāhgīr)}]{\hi{राहगीर} (rāhgīr) — a passer-by}

\label{lesson:HI-C403-rahgir}



\begin{quote}
Before the new one: say the Hindi for to rest a while, then the Hindi for a stand.
\end{quote}

\subsection*{You'll want to know: \hi{राहगीर}}
\textbf{\hi{राहगीर}} — \emph{rāhgīr} — \textquotedblleft{}a passer-by\textquotedblright{}.

\begin{quote}
\hi{एक} \hi{राहगीर} \hi{ने} \hi{हमें} \hi{रास्ता} \hi{बताया।}
\end{quote}

\begin{quote}
\emph{ek rāhgīr ne hamẽ rāstā batāyā.}
\end{quote}

\begin{quote}
A passer-by showed us the way.
\end{quote}

Masculine or feminine, by the person meant.

\begin{cousinweb}
\textbf{\hi{राहगीर}} (\emph{rāhgīr}) is usually traced to Persian \emph{rāhgīr}, \textquotedblleft{}one who takes the road\textquotedblright{}, from \emph{rāh}, \textquotedblleft{}road\textquotedblright{}: anyone walking past whom you might ask the way.
\end{cousinweb}

\begin{grammarlens}[title={What you've built}]
One more word for getting around.
\end{grammarlens}

\subsection*{Your turn}
\begin{itemize}
\raggedright
  \item \emph{Say it:} \emph{rāhgīr}
  \item \emph{Say it:} \emph{rāhgīr}, once more
  \item \emph{Say it:} \emph{ek rāhgīr ne hamẽ rāstā batāyā.} — A passer-by showed us the way.
  \item \emph{Recall:} say the Hindi for to rest a while, then the Hindi for a stand, then say \emph{rāhgīr} again
\end{itemize}

\begin{tcolorbox}[breakable,colback=teal!4,colframe=teal!35!black,title={Before you move on}]
What does \hi{राहगीर} mean? (\textquotedblleft{}A passer-by\textquotedblright{}.) Say it once more.
\end{tcolorbox}
\section[\texorpdfstring{\hi{फ़ुटपाथ} (fuṭpāth)}{फ़ुटपाथ (fuṭpāth)}]{\hi{फ़ुटपाथ} (fuṭpāth) — a pavement, a footpath}

\label{lesson:HI-C403-futpath}



\begin{quote}
Before the new one: say the Hindi for a stand, then the Hindi for a passer-by.
\end{quote}

\subsection*{You'll want to know: \hi{फ़ुटपाथ}}
\textbf{\hi{फ़ुटपाथ}} — \emph{fuṭpāth} — \textquotedblleft{}a pavement, a footpath\textquotedblright{}.

\begin{quote}
\hi{फ़ुटपाथ} \hi{पर} \hi{बहुत} \hi{भीड़} \hi{थी।}
\end{quote}

\begin{quote}
\emph{fuṭpāth par bahut bhīṛ thī.}
\end{quote}

\begin{quote}
The pavement was very crowded.
\end{quote}

Masculine.

\begin{cousinweb}
\textbf{\hi{फ़ुटपाथ}} (\emph{fuṭpāth}) is borrowed from English \emph{footpath}, in India the raised pavement beside a city road.
\end{cousinweb}

\begin{grammarlens}[title={What you've built}]
One more word for getting around.
\end{grammarlens}

\subsection*{Your turn}
\begin{itemize}
\raggedright
  \item \emph{Say it:} \emph{fuṭpāth}
  \item \emph{Say it:} \emph{fuṭpāth}, once more
  \item \emph{Say it:} \emph{fuṭpāth par bahut bhīṛ thī.} — The pavement was very crowded.
  \item \emph{Recall:} say the Hindi for a stand, then the Hindi for a passer-by, then say \emph{fuṭpāth} again
\end{itemize}

\begin{tcolorbox}[breakable,colback=teal!4,colframe=teal!35!black,title={Before you move on}]
What does \hi{फ़ुटपाथ} mean? (\textquotedblleft{}A pavement, a footpath\textquotedblright{}.) Say it once more.
\end{tcolorbox}
\section[\texorpdfstring{\hi{टहलना} (ṭahalnā)}{टहलना (ṭahalnā)}]{\hi{टहलना} (ṭahalnā) — to take a stroll; to walk up and down}

\label{lesson:HI-C403-tahalna}



\begin{quote}
Before the new one: say the Hindi for a passer-by, then the Hindi for a pavement.
\end{quote}

\subsection*{You'll want to know: \hi{टहलना}}
\textbf{\hi{टहलना}} — \emph{ṭahalnā} — \textquotedblleft{}to take a stroll; to walk up and down\textquotedblright{}.

\begin{quote}
\hi{हम} \hi{शाम} \hi{को} \hi{पार्क} \hi{में} \hi{टहलते} \hi{हैं।}
\end{quote}

\begin{quote}
\emph{ham śām ko pārk meṁ ṭahalte haiṁ.}
\end{quote}

\begin{quote}
We stroll in the park in the evening.
\end{quote}

Past (he …): \hi{वह} \hi{नदी} \hi{के} \hi{किनारे} \hi{टहला}, \emph{vah nadī ke kināre ṭahlā}.

\begin{cousinweb}
\textbf{\hi{टहलना}} (\emph{ṭahalnā}) goes back to \emph{*ṭahall}, \textquotedblleft{}walk up and down\textquotedblright{}, a form CDIAL reconstructs (CDIAL 5453), which lists Hindi \emph{ṭahalnā}, \textquotedblleft{}to take a walk\textquotedblright{}. \hi{घूमना} wanders; \hi{टहलना} strolls for the pleasure of it.
\end{cousinweb}

\begin{grammarlens}[title={What you've built}]
One more word for getting around.
\end{grammarlens}

\subsection*{Your turn}
\begin{itemize}
\raggedright
  \item \emph{Say it:} \emph{ṭahalnā}
  \item \emph{Say it:} \emph{ṭahalnā}, once more
  \item \emph{Say it:} \emph{ham śām ko pārk meṁ ṭahalte haiṁ.} — We stroll in the park in the evening.
  \item \emph{Recall:} say the Hindi for a passer-by, then the Hindi for a pavement, then say \emph{ṭahalnā} again
\end{itemize}

\begin{tcolorbox}[breakable,colback=teal!4,colframe=teal!35!black,title={Before you move on}]
What does \hi{टहलना} mean? (\textquotedblleft{}To take a stroll; to walk up and down\textquotedblright{}.) Say it once more.
\end{tcolorbox}
\section[\texorpdfstring{\hi{हाँफना} (hā̃phnā)}{हाँफना (hā̃phnā)}]{\hi{हाँफना} (hā̃phnā) — to pant}

\label{lesson:HI-C403-haphna}



\begin{quote}
Before the new one: say the Hindi for a pavement, then the Hindi for to take a stroll.
\end{quote}

\subsection*{You'll want to know: \hi{हाँफना}}
\textbf{\hi{हाँफना}} — \emph{hā̃phnā} — \textquotedblleft{}to pant\textquotedblright{}.

\begin{quote}
\hi{सीढ़ियाँ} \hi{चढ़कर} \hi{मैं} \hi{हाँफ} \hi{रहा} \hi{था।}
\end{quote}

\begin{quote}
\emph{sīṛhiyā̃ caṛhkar maiṁ hā̃ph rahā thā.}
\end{quote}

\begin{quote}
I was panting after climbing the stairs.
\end{quote}

Past (he …): \hi{वह} \hi{बहुत} \hi{हाँफा}, \emph{vah bahut hā̃phā}.

\begin{cousinweb}
\textbf{\hi{हाँफना}} (\emph{hā̃phnā}) goes back to \emph{*hamph}, \textquotedblleft{}pant\textquotedblright{}, a form CDIAL reconstructs (CDIAL 13973), which lists Hindi \emph{hā̃phnā}; the word itself sounds like the breathing it names.
\end{cousinweb}

\begin{grammarlens}[title={What you've built}]
One more word for getting around.
\end{grammarlens}

\subsection*{Your turn}
\begin{itemize}
\raggedright
  \item \emph{Say it:} \emph{hā̃phnā}
  \item \emph{Say it:} \emph{hā̃phnā}, once more
  \item \emph{Say it:} \emph{sīṛhiyā̃ caṛhkar maiṁ hā̃ph rahā thā.} — I was panting after climbing the stairs.
  \item \emph{Recall:} say the Hindi for a pavement, then the Hindi for to take a stroll, then say \emph{hā̃phnā} again
\end{itemize}

\begin{tcolorbox}[breakable,colback=teal!4,colframe=teal!35!black,title={Before you move on}]
What does \hi{हाँफना} mean? (\textquotedblleft{}To pant\textquotedblright{}.) Say it once more.
\end{tcolorbox}
